% ==============================================================================
% frontmatter/preface.tex — คำนำ
% ==============================================================================
\chapter*{คำนำ}
\addcontentsline{toc}{chapter}{คำนำ}

การประยุกต์ใช้เทคโนโลยีปัญญาประดิษฐ์ (Artificial Intelligence: AI) และระบบการเกษตรแม่นยำ (Precision Agriculture) กำลังมีบทบาทสำคัญอย่างยิ่งต่อการยกระดับภาคการเกษตรของประเทศไทย โดยเฉพาะในพื้นที่ภาคตะวันออก เช่น จังหวัดจันทบุรี ซึ่งเป็นศูนย์กลางการผลิตไม้ผลเศรษฐกิจมูลค่าสูง เช่น ทุเรียน มังคุด และพืชไร่เศรษฐกิจสำคัญ เช่น ข้าว ข้าวโพด และมันสำปะหลัง การพัฒนาขีดความสามารถของเกษตรกร นักวิจัย และผู้ประกอบการนวัตกรรมให้สามารถเข้าใจสถาปัตยกรรมข้อมูล และสามารถฝึกฝนโมเดล AI เฉพาะทางด้วยตนเอง จึงเป็นกุญแจสำคัญสู่ความสำเร็จที่ยั่งยืน

คู่มือฉบับนี้จัดทำขึ้นเพื่อใช้ประกอบการอบรมเชิงปฏิบัติการ (Workshop Hands-on) โครงการ Smart Soil \& Plant Environmental AI Analyzer โดยเน้นการวางรากฐานสถาปัตยกรรมชุดข้อมูลอย่างเป็นระบบ แบ่งออกเป็นสองเสาหลัก ได้แก่ การวิเคราะห์ดิน (Soil Domain) และการวินิจฉัยสุขภาพพืช (Plant Domain) พร้อมทั้งแยกแยะระหว่างชุดข้อมูลภาพถ่ายจริงจากงานวิจัยระดับโลก (Real Benchmark Data) และชุดข้อมูลจำลอง (Synthetic Simulation Data) เพื่อความเป็นเลิศทางวิชาการและการทดสอบที่โปร่งใส

ผู้เขียนหวังเป็นอย่างยิ่งว่า คู่มือฉบับนี้จะเป็นประโยชน์ต่อคณาจารย์ นักศึกษา นักวิจัย ตลอดจนผู้สนใจทั่วไปในการนำเทคโนโลยีปัญญาประดิษฐ์ การประมวลผลภาพดิจิทัล และฟิสิกส์เกษตรดิจิทัลไปต่อยอดสร้างสรรค์นวัตกรรมที่เป็นรูปธรรมต่อไป

\vspace{0.3in}
\begin{flushright}
\textbf{ผศ.ดร.ชีวะ ทัศนา}\\
มหาวิทยาลัยราชภัฏรำไพพรรณี\\
ตุลาคม 2569
\end{flushright}

\cleardoublepage
